\chapter{Frontend-Implementierung}
\label{kap:frontend}
\verfasser{\autorDG}

\section{Aufbau des React-Clients}

Der Client im Verzeichnis \texttt{client/} ist mit Vite als Build-Werkzeug und TypeScript umgesetzt. Für das Layout wird Tailwind CSS verwendet, für Eingabeelemente wie Schaltflächen, Checkboxen und Radiobuttons Komponenten auf Basis von Radix UI. Kurzzeitige Hinweise (Toasts) werden mit der Bibliothek \texttt{sonner} angezeigt.

Die Oberfläche besteht aus drei Seiten und mehreren Komponenten:

\begin{description}
  \item[\texttt{Auth}] Anmeldung und Registrierung
  \item[\texttt{MainMenu}] Liste der Lobbys, Dialoge zum Erstellen einer Lobby (\texttt{CreateLobbyModal}) und zur Passworteingabe (\texttt{JoinPasswordModal})
  \item[\texttt{LobbyScreen}] Spielerliste, Bereit-Status, Einladen und Lobby-Chat
  \item[\texttt{UserProfile}] Menü oben rechts mit Freundesliste, Freundschaftsanfragen und Einladungen
  \item[\texttt{PrivateChat}] Panel unten links für private Nachrichten
\end{description}

Der Client verwendet keinen Router. Welche Seite angezeigt wird, bestimmt das Feld \texttt{currentPage} im Store, und \texttt{App.tsx} rendert anhand dieses Werts die passende Seitenkomponente. Für einen Seitenwechsel setzt der Client nur diesen Wert neu, etwa wenn der Server den Beitritt zu einer Lobby bestätigt. React tauscht daraufhin die angezeigte Komponente aus, ohne dass der Browser eine neue Seite lädt.

\texttt{Auth} sendet Benutzername und Passwort mit \texttt{fetch} an \texttt{/login} bzw.\ \texttt{/register}. Die Pfade sind relativ, weil Client und Server denselben Ursprung haben. Ist die Antwort erfolgreich, ruft \texttt{Auth} \texttt{connect()} auf. Die Tokens selbst kommen im Client-Code nicht vor: Der Browser speichert die Cookies aus der Antwort und sendet das Access-Token beim WebSocket-Handshake automatisch mit. Beim nächsten Aufruf der Seite versucht der Client sofort, die Verbindung aufzubauen. Ist das Access-Token abgelaufen, das Refresh-Token aber noch gültig, erneuert der Client die Sitzung automatisch (Abschnitt~\ref{sec:wsauth}). Der Nutzer gelangt so ohne erneute Anmeldung ins Hauptmenü.

\section{Verbindungsverwaltung im Hook \texttt{useWebSocket}}

Die gesamte Kommunikation mit dem Server ist im Hook \texttt{useWebSocket} gekapselt. Er wird einmal in \texttt{App.tsx} aufgerufen und liefert Funktionen für alle Aktionen, etwa \texttt{createLobby}, \texttt{joinLobby}, \texttt{sendLobbyChatMessage} und \texttt{sendPrivateMessage}. \texttt{App.tsx} reicht diese Funktionen als Props an die Seiten weiter. Jede dieser Funktionen erzeugt eine Nachricht mit dem passenden \texttt{type} und übergibt sie an \texttt{sendMessage()}. Diese sendet nur, wenn die Verbindung offen ist, und zeigt andernfalls einen Hinweis an.

Die Verbindung selbst liegt in einer Ref (\texttt{useRef}), weil ihr Austausch kein erneutes Rendern auslösen soll. Listing~\ref{lst:connect} zeigt die Funktion \texttt{connect()}.

\begin{lstlisting}[language=JavaScript, caption={Verbindungsaufbau in \texttt{useWebSocket.tsx} (gekürzt)}, label=lst:connect]
// wss = WebSocket über TLS, analog zu https
const wsUrl = `${location.protocol === 'https:' ? 'wss' : 'ws'}://${location.host}/ws`;

const connect = (afterRefresh = false) => {
  if (ws.current && (ws.current.readyState === WebSocket.CONNECTING ||
                     ws.current.readyState === WebSocket.OPEN)) return;

  const socket = new WebSocket(wsUrl); // Access-Token-Cookie sendet der Browser mit
  ws.current = socket;
  let opened = false;

  socket.onopen = () => {
    if (ws.current !== socket) return;
    opened = true; // der Server sendet "welcome" von sich aus
  };

  socket.onmessage = (event) => {
    if (ws.current !== socket) return;
    const data = JSON.parse(event.data);
    switch (data.type) { /* ein case je Nachrichtentyp */ }
  };

  socket.onclose = (event) => {
    if (ws.current !== socket) return;
    ws.current = null;
    if (!opened) {                              // Handshake abgelehnt
      if (afterRefresh) { resetToAuth(); return; }
      refreshSession().then((ok) => ok ? connect(true) : resetToAuth());
    } else if (event.code === 4000) resetToAuth();          // abgemeldet
    else if (event.code === 1008) { resetToAuth(); /* Hinweis anzeigen */ }
  };
};
\end{lstlisting}

Die Variable \texttt{opened} unterscheidet, warum eine Verbindung endet. Lehnt der Server den Handshake ab, weil kein gültiges Access-Token vorliegt, wird \texttt{onclose} aufgerufen, ohne dass \texttt{onopen} je ausgeführt wurde. Der Client erneuert dann einmal die Sitzung über \texttt{/auth/refresh} und verbindet sich erneut. Scheitert auch das, zeigt er die Anmeldeseite. Schließt der Server dagegen eine bereits offene Verbindung mit dem Code \texttt{4000} (Abmeldung) oder \texttt{1008} (Anmeldung an einem anderen Gerät), kehrt der Client direkt zur Anmeldeseite zurück. Die Prüfung \texttt{ws.current !== socket} am Anfang jedes Handlers sorgt dafür, dass Ereignisse eines bereits ersetzten Sockets ignoriert werden.

\section{Zustand-Store als Brücke zwischen Socket und Oberfläche}

Der Store in \texttt{store.ts} enthält alle Daten, die die Oberfläche benötigt: den angemeldeten Spieler, die Lobby-Liste, die aktuelle Lobby mit Chat und Tipp-Anzeige, Freunde, Freundschaftsanfragen, Einladungen, die privaten Chats und die aktuelle Seite. Zu jedem Feld gibt es Funktionen, die es verändern. Sie erzeugen stets neue Objekte bzw.\ Arrays, anstatt bestehende zu verändern. Nur so erkennen Zustand und React, dass sich etwas geändert hat.

Der \texttt{onmessage}-Handler ruft die Funktionen des Stores über \texttt{useStore.getState()} auf (Listing~\ref{lst:store}). Da \texttt{getState()} immer den aktuellen Zustand liefert, spielt es keine Rolle, dass der Handler nur einmal beim Verbindungsaufbau registriert wird. Die Komponenten lesen den Store über Selektoren. \texttt{MainMenu} abonniert zum Beispiel nur \texttt{broadcastedLobbies} und wird deshalb nur neu gerendert, wenn sich die Lobby-Liste ändert. Abbildung~\ref{fig:datenfluss} zeigt den Weg einer Nachricht vom Socket bis zur Anzeige und den umgekehrten Weg einer Nutzeraktion.

\begin{lstlisting}[language=JavaScript, caption={Verarbeitung einer Nachricht und Lesen im Store (gekürzt)}, label=lst:store]
// useWebSocket.tsx: eingehende Nachricht in den Store schreiben
case MessageTypes.ResponseLobbyList:
  useStore.getState().setBroadcastedLobbies(data.lobbies);
  break;

// store.ts: Funktion im Store
setBroadcastedLobbies: (lobbies) => set({ broadcastedLobbies: lobbies }),

// MainMenu.tsx: Komponente liest nur die Lobby-Liste
const lobbies = useStore((state) => state.broadcastedLobbies);
\end{lstlisting}

\begin{figure}[ht]
\centering
\resizebox{\linewidth}{!}{%
\begin{tikzpicture}[font=\sffamily\small, >=Stealth,
  box/.style={draw, rounded corners=2pt, fill=white, minimum height=1.1cm, align=center, minimum width=2.8cm},
  lbl/.style={font=\sffamily\scriptsize, align=center, inner sep=2pt}]
  \node[box] (srv)   at (0,0)    {Server};
  \node[box] (hook)  at (5,0)    {Hook\\\texttt{onmessage}};
  \node[box] (store) at (10,0)   {Zustand-Store};
  \node[box] (comp)  at (15,0)   {Komponente};
  % eingehend (oben)
  \draw[->, thick] ([yshift=0.25cm]srv.east) -- node[lbl, above, yshift=0.35cm] {JSON-Nachricht} ([yshift=0.25cm]hook.west);
  \draw[->, thick] ([yshift=0.25cm]hook.east) -- node[lbl, above, yshift=0.35cm] {\texttt{getState().setX()}} ([yshift=0.25cm]store.west);
  \draw[->, thick] ([yshift=0.25cm]store.east) -- node[lbl, above, yshift=0.35cm] {Selektor, Re-Rendering} ([yshift=0.25cm]comp.west);
  % ausgehend (unten)
  \draw[->] ([yshift=-0.25cm]hook.west) -- node[lbl, below, yshift=-0.35cm] {JSON-Nachricht über \texttt{sendMessage()}} ([yshift=-0.25cm]srv.east);
  \draw[->] (comp.south) -- ++(0,-1.2) -| node[lbl, pos=0.25, below] {Aufruf einer Hook-Funktion über Props (\texttt{joinLobby}, \texttt{sendLobbyChatMessage}, \ldots)} (hook.south);
\end{tikzpicture}}
\caption{Datenfluss im Client: eingehende Nachrichten (oben) und Nutzeraktionen (unten)}
\label{fig:datenfluss}
\end{figure}

Bei manchen Nachrichten leitet der Client aus dem Unterschied zwischen altem und neuem Zustand zusätzliche Informationen ab. Bei \texttt{lobby\_updated} vergleicht er die Spielerliste der bisherigen Lobby mit der neuen und zeigt für jeden hinzugekommenen bzw.\ fehlenden Spieler einen Hinweis an (\enquote{lena joined the lobby}).

Für die privaten Chats speichert der Store je Freund die Nachrichten, die Zahl der ungelesenen Nachrichten, den Zeitpunkt der letzten Aktivität und ob der Freund gerade tippt. Welche Nachricht zuletzt gelesen wurde, wird zusätzlich im \texttt{localStorage} abgelegt. Der Schlüssel enthält die IDs beider Spieler, da sich mehrere Konten einen Browser teilen können. Dadurch stimmen die Zähler auch nach einem Neuladen.

\section{Benutzeroberfläche}

Die Seiten zeigen ausschließlich Daten aus dem Store an und rufen für Aktionen die Funktionen aus \texttt{useWebSocket} auf. Ein Klick auf eine Lobby im Hauptmenü ruft \texttt{joinLobby} auf. Die Lobby-Ansicht erscheint aber erst, wenn der Server mit \texttt{join\_lobby\_successful} antwortet und der Store die Seite wechselt. Die Oberfläche zeigt damit immer den Zustand, den der Server bestätigt hat.

Abbildung~\ref{fig:freunde} zeigt das Profilmenü mit der Freundesliste. Der grüne Punkt neben jedem Namen gibt den Online-Status an. Das Feld darüber dient zum Senden von Freundschaftsanfragen, darunter stehen offene Anfragen mit Schaltflächen zum Annehmen und Ablehnen. Hat der Spieler offene Anfragen oder Einladungen, zeigt das Profilsymbol einen roten Zähler. Die Lobby-Ansicht und der private Chat werden in Kapitel~\ref{kap:umsetzung} bei den jeweiligen Anforderungen gezeigt.

\begin{figure}[ht]
  \centering
  \screenshot[0.8\linewidth]{freunde.png}
  \caption{Profilmenü mit Freundesliste und Online-Status}
  \label{fig:freunde}
\end{figure}
